\chapter{闭区间上的连续函数：习题解答}

\section*{A组}
\addcontentsline{toc}{section}{A组}

\begin{solution}{18.1}
对 \((0,1)\) 取 \(U_n=(1/n,1)\)；对 \([0,\infty)\) 取 \(U_n=(-1,n)\)；对 \(\R\) 取 \(U_n=(-n,n)\)。三族都分别覆盖目标集。任意有限子族只需取最大指标 \(N\)，其并分别仍漏掉靠近 \(0\) 的正点、大于等于 \(N\) 的点或绝对值不小于 \(N\) 的点。
\end{solution}

\begin{solution}{18.2}
设 \(\mathcal U\) 是 \([a,b]\) 的一个开覆盖，并令
\[
 E=\{x\in[a,b]:[a,x]\text{ 可由 }\mathcal U\text{ 中有限多个集合覆盖}\}.
\]
取 \(U_a\in\mathcal U\) 使 \(a\in U_a\)，便知 \(a\in E\)。集合 \(E\) 非空且以 \(b\) 为上界，故
\[
 c=\sup E
\]
存在。

反设 \(c<b\)。取 \(V\in\mathcal U\) 使 \(c\in V\)，再取 \(\rho>0\)，使
\[
 (c-\rho,c+\rho)\cap[a,b]\subset V,
 \qquad c+\rho/2\le b.
\]
由上确界的逼近性质，可选 \(x\in E\) 满足 \(x>c-\rho/2\)。用有限多个覆盖成员覆盖 \([a,x]\)，再添上 \(V\)，便覆盖了 \([a,c+\rho/2]\)。于是 \(c+\rho/2\in E\)，与 \(c\) 是上界矛盾。因此 \(c=b\)。

等式 \(\sup E=b\) 还没有说明 \(b\in E\)。为此取 \(U_b\in\mathcal U\) 使 \(b\in U_b\)，并取 \(\eta>0\) 使 \((b-\eta,b]\subset U_b\)。由 \(\sup E=b\)，可选 \(x\in E\) 使 \(b-\eta<x\le b\)。若 \(x=b\)，结论已经成立；若 \(x<b\)，覆盖 \([a,x]\) 的有限多个集合再加 \(U_b\) 覆盖整个 \([a,b]\)。故 \(b\in E\)，有限覆盖得证。
\end{solution}

\begin{solution}{18.3}
上确界不必属于原集合。例如 \(\sup(0,1)=1\notin(0,1)\)。因此 \(c=b\) 只说明可有限覆盖的左端区间能任意接近 \(b\)，尚未说明整个 \([a,b]\) 可覆盖。利用含 \(b\) 的开集向左覆盖一小段，再与某个足够靠近 \(b\) 的已覆盖区间拼接，才得到 \(b\in E\)。
\end{solution}

\begin{solution}{18.4}
若 \(K\) 闭且有界，取 \(K\subset[-M,M]\)。给定 \(K\) 的开覆盖，添入开集 \(\R\setminus K\)，便得到 \([-M,M]\) 的开覆盖。闭区间定理给出有限子覆盖。删去 \(\R\setminus K\) 后，余下原覆盖成员仍覆盖 \(K\)。
\end{solution}

\begin{solution}{18.5}
若 \(K\) 无界，则 \(\{(-n,n):n\in\Npos\}\) 是 \(K\) 的开覆盖。任意有限子族的并为某个 \((-N,N)\)，不能覆盖无界的 \(K\)，与有限覆盖性质矛盾。
\end{solution}

\begin{solution}{18.6}
若 \(K\) 不闭，由闭集的序列刻画存在 \(x_j\in K\) 与 \(a\notin K\)，使 \(x_j\to a\)。令
\[
 U_n=\{x\in\R:\abs{x-a}>1/n\},\qquad n\in\Npos.
\]
因 \(a\notin K\)，这些开集覆盖 \(K\)。有限子族的并等于其中指标最大的一个 \(U_N\)，但 \(a\) 是 \(K\) 的聚点，存在 \(x\in K\) 满足 \(\abs{x-a}<1/N\)，故 \(x\notin U_N\)。矛盾。
\end{solution}

\begin{solution}{18.7}
覆盖法：每点连续性给出邻域 \(U_x\)，使其中函数值与 \(f(x)\) 相差小于 \(1\)。取有限子覆盖 \(U_{x_1},\ldots,U_{x_m}\)，以
\[
 1+\max_j\abs{f(x_j)}
\]
为全局界。

序列法：若无界，选 \(x_n\) 使 \(\abs{f(x_n)}>n\)。闭区间中的点列有收敛子列，其极限仍在闭区间；连续性使像子列收敛而有界，与构造矛盾。
\end{solution}

\begin{solution}{18.8}
设 \(K\subset\R\) 非空且紧，\(f:K\to\R\) 连续。其值集非空且有界，故可令 \(M=\sup f(K)\)。选 \(x_n\in K\) 使
\[
 M-1/n<f(x_n)\le M.
\]
紧集的序列刻画给出 \(x_{n_k}\to x_*\in K\)。连续性与夹逼给出
\[
 f(x_*)=\lim_k f(x_{n_k})=M.
\]
最大值取得。对 \(-f\) 重复，得到最小值。
\end{solution}

\begin{solution}{18.9}
删去闭性：\(f(x)=x\) 在 \((0,1)\) 上不取极值。删去有界性：\(f(x)=x\) 在闭集 \(\R\) 上无界。删去连续性：在 \([0,1]\) 上令
\[
 f(0)=0,\qquad f(x)=1/x\quad(x>0),
\]
则无界，因而没有最大值。
\end{solution}

\section*{B组}
\addcontentsline{toc}{section}{B组}

\begin{solution}{18.10}
反设区间 \(I\) 可写成两个不交的非空相对开集 \(A,B\) 之并。取 \(a\in A\)、\(b\in B\)，交换二者及集合名称后可设 \(a<b\)。令
\[
 E=A\cap[a,b],\qquad c=\sup E.
\]
集合 \(E\) 非空且有上界，而区间性质给出 \([a,b]\subset I\)，故 \(c\in I=A\cup B\)。

若 \(c\in A\)，由 \(A\) 在 \(I\) 中相对开，存在 \(\eta>0\) 使
\[
 (c-\eta,c+\eta)\cap I\subset A.
\]
又因 \(b\in B\)，有 \(c<b\)。于是可在 \((c,\min\{c+\eta,b\})\) 中取一点；该点属于 \(E\) 且大于 \(c\)，与 \(c\) 是上界矛盾。

若 \(c\in B\)，相对开性给出 \(c\) 在 \(I\) 中的一个邻域，其中没有 \(A\) 中点。另一方面，由上确界的逼近性质，可在这个邻域内选到 \(x\in E\) 且 \(x\le c\)；因 \(c\notin A\)，实际有 \(x<c\)。这同样矛盾。故所设分离不存在，任意区间都连通。
\end{solution}

\begin{solution}{18.11}
若连通集 \(E\) 不具区间性质，存在 \(x<z<y\)，其中 \(x,y\in E\)、\(z\notin E\)。则
\[
 E\cap(-\infty,z),\qquad E\cap(z,\infty)
\]
是两个非空、不交的相对开集，其并为 \(E\)，构成分离，矛盾。
\end{solution}

\begin{solution}{18.12}
取无理数 \(\alpha\)。则
\[
 \Q\cap(-\infty,\alpha),\qquad \Q\cap(\alpha,\infty)
\]
是 \(\Q\) 的非空相对开分离。对无理数集，取任意有理数 \(r\)，用 \(r\) 左右两个半轴与无理数集相交即可分离。两集合虽都在实线中稠密，仍不连通。
\end{solution}

\begin{solution}{18.13}
设 \(f:[a,b]\to\R\) 连续，且 \(f(a)<0<f(b)\)。令
\[
 E=\{x\in[a,b]:f(x)<0\},
 \qquad c=\sup E.
\]
由于 \(a\in E\) 而 \(E\subset[a,b]\)，数 \(c\) 存在且属于 \([a,b]\)。

若 \(f(c)<0\)，在 \(c\) 处的连续性给出 \(\eta>0\)，使 \([a,b]\) 中满足 \(\abs{x-c}<\eta\) 的点仍有 \(f(x)<0\)。此时 \(c\ne b\)，因而可在 \(c\) 右侧取这样的点；它属于 \(E\) 且大于 \(c\)，与上界性矛盾。

若 \(f(c)>0\)，连续性给出 \(c\) 的一个相对邻域，其中函数值都为正。此时 \(c\notin E\)，而上确界的逼近性质要求该邻域在 \(c\) 左侧含有 \(E\) 中点，又得到矛盾。前一次连续性把负号延续到 \(c\) 的右侧，后一次把正号延续到 \(c\) 的左侧。因此只有 \(f(c)=0\)。
\end{solution}

\begin{solution}{18.14}
给定介于 \(f(x_0),f(x_1)\) 之间的 \(y\)，令 \(g(x)=f(x)-y\)。若 \(y\) 等于某端点值，直接取该端点；否则 \(g\) 在两端严格异号。对 \(g\) 应用零点定理得到 \(g(c)=0\)，即 \(f(c)=y\)。
\end{solution}

\begin{solution}{18.15}
任取 \(u,v\in f(I)\)，选 \(x_0,x_1\in I\) 使 \(f(x_0)=u,f(x_1)=v\)。对每个介于 \(u,v\) 之间的 \(y\)，介值定理给出位于 \(x_0,x_1\) 之间的 \(c\in I\)，使 \(f(c)=y\)。所以 \(f(I)\) 具有区间性质。
\end{solution}

\begin{solution}{18.16}
若存在 \(x_0,x_1\) 使整数值 \(f(x_0)<f(x_1)\)，则两整数之间可取非整数 \(y\)。介值定理会给出 \(c\) 使 \(f(c)=y\)，与值域包含于 \(\Z\) 矛盾。因此任意两点函数值相同，函数为常值。
\end{solution}

\begin{solution}{18.17}
令 \(g(x)=f(x)-x\)。函数 \(g\) 连续，且
\[
 g(a)=f(a)-a\ge0,\qquad
 g(b)=f(b)-b\le0.
\]
若某端点值为零即完成；否则两端异号，由零点定理存在 \(c\) 使 \(g(c)=0\)，即 \(f(c)=c\)。
\end{solution}

\begin{solution}{18.18}
设多项式次数为奇数，首项为 \(a_{2m+1}x^{2m+1}\)。低次项与首项之比在 \(\abs x\to\infty\) 时趋零，因此当 \(\abs x\) 足够大时，多项式符号与首项相同。正负无穷方向上首项符号相反，故可取 \(A>0\) 使 \(p(-A)p(A)<0\)。由连续性和零点定理，\(p\) 在 \((-A,A)\) 内有实根。
\end{solution}

\section*{C组}
\addcontentsline{toc}{section}{C组}

\begin{solution}{18.19}
给定 \(f(K)\) 的开覆盖 \(\{V_\lambda\}\)。连续性使 \(f^{-1}(V_\lambda)\) 在 \(K\) 中相对开，并覆盖 \(K\)。写成 \(K\cap U_\lambda\) 后，由 \(K\) 的有限覆盖性质取有限多个原像覆盖 \(K\)。对应的有限多个 \(V_\lambda\) 覆盖 \(f(K)\)。
\end{solution}

\begin{solution}{18.20}
若 \(f(E)=A\cup B\) 可由两个非空、不交的相对开集分离，则
\[
 f^{-1}(A),\qquad f^{-1}(B)
\]
是 \(E\) 的两个非空、不交、相对开的集合，且并为 \(E\)。这与 \(E\) 连通矛盾，故连续像连通。
\end{solution}

\begin{solution}{18.21}
设 \(S\subset[a,b]\) 无限。把区间二分，至少一个半区间含 \(S\) 的无限多个点；递归选择这种闭半区间，得到套区间 \(I_n\)，长度趋零且每个都含 \(S\) 的无限点。套区间定理给出唯一公共点 \(c\)。每个 \(c\) 的邻域包含某个 \(I_n\)，其中有异于 \(c\) 的 \(S\) 中点，故 \(c\) 是聚点。

再用已经得到的聚点性质证明有限覆盖定理。若 \([a,b]\) 有一个无有限子覆盖的开覆盖，必有 \(a<b\)。逐次二分并保留不能有限覆盖的闭半区间，得到 \(J_n\)，其长度趋零。记中心为 \(c_n\)。后继半区间的中心不等于先前中心，因为先前中心是后继区间的端点或在其外；故中心组成无限集。第一问给出这个集合的聚点 \(c\in[a,b]\)。

固定 \(m\)，除有限个中心外其余都在闭集 \(J_m\) 中，所以 \(c\in J_m\)：若不在，可用 \(J_m\) 的开补集取得避开所有后续中心的邻域，再缩小以避开其余异于 \(c\) 的有限个中心，违背聚点性。因此 \(c\) 属于全部 \(J_m\)。取覆盖中含 \(c\) 的开集 \(U\) 与 \(r>0\)，使 \((c-r,c+r)\subset U\)。当 \(J_m\) 的长度小于 \(r\) 时，因 \(c\in J_m\)，便有 \(J_m\subset U\)，与它不能有限覆盖矛盾。此段使用第一问的聚点结论，不再直接调用区间套的公共点结论。
\end{solution}

\begin{solution}{18.22}
设 \(f:K\to Y\subset\R\) 是连续双射，\(K\) 紧致。记 \(g=f^{-1}\)。若 \(C\subset K\) 在 \(K\) 中相对闭，则 \(C\) 紧；连续像 \(f(C)\) 紧，因而在实线中闭，也在 \(Y\) 中相对闭。由于 \(g^{-1}(C)=f(C)\)，逆映射满足闭集原像判据，故连续。这里研究的是原映射 \(f\) 的闭集像，而不是假设逆映射已经保持某类像。
\end{solution}

\begin{solution}{18.23}
先证闭且有界推出序列紧致。任取 \((x_n)\subset K\)，有界性与 Bolzano--Weierstrass 定理给出收敛子列；闭性保证其极限仍在 \(K\) 中。

再由序列紧致推出闭且有界。若 \(K\) 无界，可递归选择 \(x_n\in K\) 使 \(\abs{x_n}>n\)，所得数列没有收敛子列，矛盾。若 \(K\) 不闭，闭集的序列刻画给出 \(x_n\in K\) 与 \(a\notin K\)，使 \(x_n\to a\)。它的每个子列仍收敛到 \(a\)，因极限唯一，不可能有子列收敛到 \(K\) 中的点。故 \(K\) 闭且有界，再由 Heine--Borel 定理得到有限覆盖性质。这里没有枚举任意开覆盖；开覆盖可能不可数，不能写成“从逐渐增大的有限子族外选点”。

最后由有限覆盖性质推出闭且有界。若无界，开覆盖 \(\{(-n,n):n\in\Npos\}\) 没有有限子覆盖。若不闭，由闭集序列刻画可取 \(x_j\in K\) 收敛到 \(a\notin K\)，则
\[
 U_n=\{x\in\R:\abs{x-a}>1/n\},\qquad n\in\Npos,
\]
覆盖 \(K\)，但任一有限子族都会漏掉充分靠近 \(a\) 的 \(K\) 中点。因此 \(K\) 必闭且有界。

完备性在第一步通过 Bolzano--Weierstrass 定理进入，在“闭且有界推出有限覆盖性质”中通过闭区间有限覆盖定理进入；后二者先前都已由实数完备性证明。其余反例选择和闭性论证不再增加新的完备性原理。
\end{solution}

\begin{solution}{18.24}
若 \(K=\varnothing\)，任意 \(\lambda>0\) 都满足结论。以下设 \(K\ne\varnothing\)。
对每个 \(x\in K\)，选 \(U_x\in\mathcal U\) 使 \(x\in U_x\)。由 \(U_x\) 开，可取 \(r_x>0\) 使
\[
 (x-2r_x,x+2r_x)\subset U_x.
\]
区间族 \(\{(x-r_x,x+r_x):x\in K\}\) 覆盖 \(K\)，故有有限子覆盖，对应中心记为 \(x_1,\ldots,x_m\)。令
\[
 \lambda=\min_{1\le j\le m}r_{x_j}>0.
\]
任取 \(y\in K\)，有某个 \(j\) 使 \(\abs{y-x_j}<r_{x_j}\)。若 \(z\in K\cap(y-\lambda,y+\lambda)\)，则
\[
 \abs{z-x_j}
 \le\abs{z-y}+\abs{y-x_j}
 <\lambda+r_{x_j}
 \le2r_{x_j},
\]
所以 \(z\in U_{x_j}\)。于是
\[
 K\cap(y-\lambda,y+\lambda)\subset U_{x_j},
\]
证明了 \(\lambda\) 是所求 Lebesgue 数。
\end{solution}

\begin{solution}{18.25}
任取 \(x_n\in K_n\)。因 \(K_1\) 紧，存在子列 \(x_{n_k}\to x\in K_1\)。固定 \(m\)。当 \(n_k\ge m\) 时，由递减性 \(x_{n_k}\in K_{n_k}\subset K_m\)。集合 \(K_m\) 闭，故极限 \(x\in K_m\)。这对每个 \(m\) 成立，所以
\[
 x\in\bigcap_{m=1}^{\infty}K_m.
\]
\(K_1\) 的有界性已由它的紧致性保证，不需重复假设。若只要求非空闭集，取 \(K_n=[n,+\infty)\)，它们递减却没有公共点。这个例子也说明不能把序列抽取所用的紧致性只换成闭性。
\end{solution}

\begin{solution}{18.26}
设 \(K\) 与 \(f\) 如题。零点集
\[
 Z=f^{-1}(\{0\})
\]
在 \(K\) 中闭，因此是紧集；这部分对空集也成立。

再设 \(f\) 在 \(K\) 上处处非零。连续函数 \(\abs f\) 由极值定理在非空紧集 \(K\) 上取得最小值 \(m\ge0\)。若 \(m=0\)，则有 \(x_0\in K\) 使 \(\abs{f(x_0)}=0\)，与处处非零矛盾，故 \(m>0\)。
\end{solution}
